\documentclass[11pt,a4paper]{article}
\usepackage[margin=1in]{geometry}
\usepackage{hyperref}
\usepackage{enumitem}
\usepackage{graphicx}
\usepackage{xcolor}
% -----------------------------------------------------
% bvraghav-tiet-ucs503p-article.sty
% -----------------------------------------------------
% Variable: Lab Instructor
\makeatletter \newcommand{\BvrSetLabInstructor}[1]{%
  \gdef\Bvr@LabInstructor{#1}}
% default
\newcommand{\Bvr@LabInstructor}{Dr. Raghav %
  B. Venkataramaiyer}
% public accessor
\newcommand{\BvrLabInstructorValue}{\Bvr@LabInstructor}
\makeatother
% Add subtitle
\usepackage{titling}
% -- Set title bold
\pretitle{\begin{center}\bfseries\fontsize{18bp}{18bp}\selectfont}
\makeatletter
\newcommand{\BvrSubtitle}[1]{%
  \posttitle{%
    \par\end{center}
    \begin{center}\large#1\end{center}
    \vskip 0.5em}%
}
\makeatother
% Hints in the section title(s)
\newcommand{\BvrHint}[1]{%
  {\normalfont\normalsize\itshape\hfill #1}}
% -----------------------------------------------------
\title{CodeSphere AI: An AI-Augmented Software Engineering Collaboration Platform}
\BvrSetLabInstructor{Ms. Paramveer Sidhu}
\BvrSubtitle{%
  Project Proposal for UCS503P  \\
  Submitted to: \BvrLabInstructorValue \\ \bfseries
  Thapar Institute of Engineering and Technology%
}
\author{%
  Sarthak Soni, 3C11%
  \thanks{Roll No: \texttt{1024030061}, %
    Email: \texttt{ssoni1\_be24\@thapar.edu}}
  \and
  Ayush Bansal, 3C11%
  \thanks{Roll No: \texttt{1024030059}, %
    Email: \texttt{abansal9\_be24\@thapar.edu}}%
}
\date{\today}
\begin{document}
\maketitle
\section{Higher-order goal%
  \BvrHint{\ldots secondary framing}}
This project aims to improve software engineering productivity by
reducing the time developers spend understanding unfamiliar
codebases, triaging issues, and coordinating work across a team.
While the underlying theme (AI-assisted developer tooling) is
broad, the immediate engineering objective is to translate
\emph{faster, better-informed engineering decisions} into a
measurable, deployable software solution.
\section{Time-to-value %
\BvrHint{\ldots primary heuristic}}
\begin{itemize}[leftmargin=0.7cm]
    \item We will deliver meaningful increments quickly by prototyping the core repository-insight workflow and validating it against a real open-source repository within a short iteration window.
    \item We will prioritize fast feedback over perfection by using weekly iterations (and targeting CI-backed automation early) to reduce release friction.
    \item Success is defined by what can be delivered and measured in the first iteration, then improved based on user feedback.
\end{itemize}
\section{Problem Statement}
Third-year CS teams and early-career developers routinely lose time
on tasks that are procedural rather than creative: understanding
what a codebase does, finding the right file/module for a bug,
writing consistent documentation, and tracking who is doing what.
Common pain points include:
\begin{itemize}[leftmargin=0.7cm]
    \item \textbf{Onboarding friction:} new contributors take long to understand project structure and module responsibilities.
    \item \textbf{Fragmented task tracking:} tasks/issues live separately from the code, so context is repeatedly re-explained.
    \item \textbf{Manual, inconsistent code review:} reviewers spend time on style/documentation nits instead of logic and design.
    \item \textbf{Poor visibility:} team leads lack a single view of project health (open issues, stale PRs, module ownership).
\end{itemize}
\textbf{Why this matters now (contextual relevance):} Generative and
retrieval-based AI models are mature and cheap enough to embed
directly into everyday developer workflows; a lightweight
integration layer can meaningfully cut down repetitive engineering
overhead even for small teams.
\section{Proposed Solution %
\BvrHint{\ldots Software Proposition}}
\subsection{Overview}
Build a \textbf{web-based} ``CodeSphere AI'' dashboard, backed by a
\textbf{GitHub App/webhook bot}, that adds AI-assisted codebase
understanding on top of GitHub rather than replacing GitHub's own
project-management features. We deliberately do \emph{not} build a
custom Kanban/task board -- GitHub Projects already covers that --
and instead surface GitHub's existing issues in our dashboard with
an AI layer on top. Components:
\begin{itemize}[leftmargin=0.7cm]
    \item \textbf{Repository connector:} link a GitHub repository via OAuth (read-only) and index its structure, commits, and open issues.
    \item \textbf{AI codebase assistant (website):} chat interface that answers questions about the codebase (``where is X handled?'', ``summarize this module'') using retrieval over indexed source files, with citations to file/line.
    \item \textbf{Auto-triage bot (GitHub App/webhook):} on every new issue/PR, the bot posts a generated summary and a suggested label as a comment directly on GitHub -- this is the feature GitHub does not provide natively.
    \item \textbf{Team dashboard (website):} read-only view of GitHub issues/PRs enriched with staleness alerts and module ownership, pulled live from the GitHub API (not a separate board).
\end{itemize}
\subsection{Core workflow}
\begin{enumerate}[leftmargin=0.7cm]
    \item Team lead connects a GitHub repository; system indexes files, commit history, and open issues.
    \item Contributor asks the AI assistant a question about the codebase (on the website) and receives an answer grounded in retrieved source snippets, with file/line references.
    \item When a new issue/PR is opened on GitHub, the webhook bot auto-comments with a summary and suggested label within seconds.
    \item Team lead views the dashboard for backlog health, stale items, and module ownership, sourced live from GitHub -- no data duplicated into a separate board.
\end{enumerate}
\subsection{Operational constraints for an educational, third-year setting}
\begin{itemize}[leftmargin=0.7cm]
    \item Keep the solution usable with a single GitHub OAuth connection and no heavyweight infrastructure.
    \item Avoid dependence on training custom ML models; focus on retrieval, orchestration, and workflow correctness using existing LLM APIs.
    \item Design for deployability using common web hosting and managed databases.
\end{itemize}
\section{Solution Approach %
\BvrHint{\ldots Engineering focus}}
This is an engineering course project, so we focus on scalable
integration, retrieval, and workflow aspects rather than novel
model research.
\subsection{Codebase understanding %
\BvrHint{\ldots non-research}}
Code intelligence will be retrieval-based, not model-training-based:
\begin{itemize}[leftmargin=0.7cm]
    \item Chunk source files by function/class boundaries and generate embeddings using an existing embedding API.
    \item Store embeddings in a vector index; retrieve top-\(k\) relevant chunks for a given question.
    \item Pass retrieved chunks plus the question to an LLM API to generate a grounded answer with citations to file paths; never claim state-of-the-art accuracy.
\end{itemize}
\subsection{Web architecture %
\BvrHint{\ldots web-based solution}}
Two cooperating pieces, sharing one backend and database:
\begin{itemize}[leftmargin=0.7cm]
    \item \textbf{Website (frontend + backend API):} responsive UI for the chat assistant and dashboard; backend handles GitHub OAuth, repository indexing jobs, and LLM orchestration.
    \item \textbf{GitHub App / webhook service:} a small service registered as a GitHub App on the repository; GitHub pushes ``issue opened'' / ``PR opened'' events to it, it calls the same LLM orchestration layer, and posts the summary/label back to GitHub via the GitHub API.
    \item \textbf{Database:} store users, repositories, indexed chunks/embeddings, and generated summaries (issues/tasks themselves stay in GitHub -- we do not duplicate that data).
\end{itemize}
\subsection{CI/CD and fast delivery enabler}
CI/CD will be used to:
\begin{itemize}[leftmargin=0.7cm]
    \item Automatically build and run tests on every change.
    \item Produce repeatable deployment artifacts to staging.
    \item Enable safe and frequent releases of incremental features.
\end{itemize}
\section{Evaluation Criterion %
\BvrHint{\ldots measurable and attributable}}
We define success using measurable metrics tied to the core workflow.
\subsection{Primary evaluation metric}
\textbf{Answer grounding rate:} percentage of AI-assistant answers
that correctly cite an existing, relevant file/line in the
indexed repository.
\begin{itemize}[leftmargin=0.7cm]
    \item Target: $\ge$ 85\% grounding rate on a held-out set of 30 test questions during pilot.
    \item Attribution: measured by comparing cited file references against a manually verified answer key.
\end{itemize}
\subsection{Secondary evaluation metrics}
\begin{itemize}[leftmargin=0.7cm]
    \item \textbf{Time-to-first-understanding:} time for a new contributor to locate the correct module for a sample task, with vs.\ without the assistant.
    \item \textbf{Issue triage accuracy:} percentage of auto-generated labels/summaries accepted without edit by the contributor.
    \item \textbf{Triage latency:} time between an issue being opened on GitHub and the bot posting its summary comment (target: under 30 seconds).
    \item \textbf{Reliability:} availability of staging/production for login and repository indexing during business hours (e.g., $\ge$ 99\% uptime in pilot).
\end{itemize}
\subsection{Pilot validation plan %
\BvrHint{\ldots high-level}}
\begin{itemize}[leftmargin=0.7cm]
    \item Recruit 2--3 small student project teams (5--8 contributors total) and index one real repository per team.
    \item Compare onboarding/triage metrics before/after within the iteration window by logging baseline estimates via short interviews or manual timing.
\end{itemize}
\section{Scalability %
\BvrHint{\ldots with foresight}}
The proposition should scale in both theoretical and practical senses.
\begin{enumerate}[leftmargin=0.7cm]
    \item \textbf{Scalable (theoretically):} The system should support growth in repository size and number of concurrent teams by:
    \begin{itemize}[leftmargin=0.7cm]
        \item decoupling components (frontend/backend/indexing worker),
        \item incremental re-indexing on new commits instead of full re-index,
        \item using stateless API design and a dedicated vector index for retrieval.
    \end{itemize}
    \item \textbf{Operationally deployable (practically):} The system should run on common web hosting:
    \begin{itemize}[leftmargin=0.7cm]
        \item managed relational database plus a managed vector store,
        \item containerized or platform-hosted deployment,
        \item CI/CD pipeline for staging/production.
    \end{itemize}
\end{enumerate}
\section{Engine availability heuristic}
A mature set of ``engines'' (tooling/services) is available to implement this as a web-based project:
\begin{itemize}[leftmargin=0.7cm]
    \item GitHub API and OAuth for repository access and authentication.
    \item Web framework for REST APIs and reactive pages.
    \item Managed relational database plus a managed vector database for embeddings.
    \item Existing LLM and embedding APIs for summarization, Q\&A, and retrieval.
    \item Test framework for unit/integration tests.
    \item CI runner and staging deployment target.
\end{itemize}
We will use these mature components to focus on workflow design, retrieval correctness, and measurement rather than inventing new infrastructure or training new models.
\section{Project scope and deliverables %
\BvrHint{\ldots iteration-friendly}}
\subsection{Initial deliverable %
\BvrHint{\ldots first iteration}}
\begin{itemize}[leftmargin=0.7cm]
    \item GitHub OAuth login + repository connection
    \item Repository indexing pipeline (chunking + embeddings) for one repository
    \item AI assistant chat with grounded, cited answers (website)
    \item GitHub App/webhook bot that auto-comments a summary on newly opened issues
    \item CI: build + backend unit tests + linting
    \item CD to staging with a single command or pipeline job
\end{itemize}
\subsection{Subsequent deliverables}
\begin{itemize}[leftmargin=0.7cm]
    \item Suggested-label posting alongside the summary (not just summary)
    \item Incremental re-indexing on new commits
    \item Team dashboard: stale-issue alerts and module ownership, pulled live from the GitHub API
    \item Multi-repository and multi-team support
\end{itemize}
\section{Risks and mitigations}
\begin{itemize}[leftmargin=0.7cm]
    \item \textbf{LLM cost/rate limits:} cache embeddings and answers; batch indexing jobs; set per-team usage quotas.
    \item \textbf{Answer hallucination:} always require retrieved-chunk citations; show ``no confident answer'' when retrieval confidence is low rather than guessing.
    \item \textbf{Repository access/privacy:} request minimal OAuth scopes; allow teams to disconnect and delete indexed data.
    \item \textbf{Evaluation measurement ambiguity:} fix the test question set and grounding-verification process before development begins.
    \item \textbf{Operational issues in pilot:} use staging early; ensure CI/CD produces repeatable deployments.
\end{itemize}
\section{Summary}
This project proposes a deployable web platform, CodeSphere AI,
that reduces onboarding and coordination overhead for student and
early-career software teams by grounding an AI assistant in a
team's own codebase and linking it to a lightweight task workflow.
It meets the course criteria by emphasizing time-to-value through
rapid iterations, implementing CI/CD to support frequent safe
releases, and using measurable evaluation criteria (especially
answer grounding rate). It remains focused on engineering
integration, retrieval correctness, and scalability readiness
rather than research-driven model novelty.
\end{document}